% !TeX root = ../../apuntes_bloque1.tex
\section{Definiciones y Esfuerzos en Secciones Inclinadas}

En un elemento estructural sometido a solicitaciones exteriores se pueden desarrollar distintos tipos de tensiones según la orientación relativa entre la fuerza aplicada y la sección resistente considerada:
\begin{itemize}
    \item \textbf{Esfuerzos normales ($\sigma$):} Se desarrollan en barras sometidas a cargas axiales aplicadas a lo largo de su eje longitudinal y actúan perpendicularmente a la sección transversal de la barra.
    \item \textbf{Esfuerzos cortantes o tangenciales ($\tau$):} Una barra soporta un esfuerzo cortante cuando las cargas transversales tienden a producir un corte o deslizamiento relativo a través de una sección $C$ comprendida entre los puntos de aplicación de las fuerzas.
\end{itemize}

\dfn{Tensión Cortante Media ($\tau$)}{Es la resistencia interna tangencial por unidad de área que desarrolla una sección resistente para equilibrar un esfuerzo de corte $P$ paralelo al plano de dicha sección:}
\begin{equation}
    \tau = \frac{P}{A}
    \label{eq:t8_tau_media}
\end{equation}
donde $P$ es la carga en cortante y $A$ es el área de la sección transversal sometida al corte. Ejemplos característicos en Ingeniería Química son los pasadores, chavetas, uniones remachadas (roblonadas) y uniones atornilladas.

\subsection{Tensiones en secciones inclinadas bajo carga axial}
Incluso cuando una barra está sometida únicamente a dos fuerzas axiales $P$ iguales y de sentido contrario a lo largo de su eje longitudinal, si en lugar de considerar la sección recta transversal $A_0$ se estudia una sección inclinada $A_\theta$ que forma un ángulo $\theta$ respecto al plano normal, aparecen simultáneamente tensiones normales y tensiones cortantes.

Cortando la barra por el plano inclinado $A_\theta$ y planteando el equilibrio del sólido libre de la izquierda, la carga axial $P$ se descompone en una componente normal $F$ perpendicular a la sección y una componente tangencial o cortante $V$ contenida en el plano de la sección:
\begin{equation}
    F = P \cos\theta, \qquad V = P \sin\theta
    \label{eq:t8_decomp_FV}
\end{equation}

\begin{figure}[H]
    \centering
    \begin{tikzpicture}[>=Latex, thick]
        % Barra izquierda con corte inclinado theta
        \draw[fill=blue!6, draw=black] (0,-0.85) -- (3.8,-0.85) -- (4.7,0.85) -- (0,0.85) -- cycle;
        \draw[dashed, gray] (0,0) -- (5.3,0);
        \draw[dashed, black!70] (4.25,-1.1) -- (4.25,1.1);
        % Angulo theta (plano inclinado forma angulo 28 deg con la vertical)
        \draw[thin, red!70!black] (4.25,0) ++(90:0.65) arc (90:62:0.65);
        \node[red!70!black, font=\small] at (4.44,0.82) {$\theta$};
        % Fuerza axial P izquierda
        \draw[->, line width=1.4pt, blue!75!black] (-1.2,0) -- (0,0) node[pos=0.35, above] {$P$};
        % Componentes F (normal al plano inclinado: dir -28 deg) y V (tangencial al plano inclinado: dir -118 deg)
        \coordinate (M) at (4.25,0);
        \draw[->, line width=1.4pt, blue!75!black] (M) -- ++(0:1.65) node[right] {$P$};
        \draw[->, line width=1.4pt, myg!80!black] (M) -- ++(-28:1.45) node[below right] {$F = P\cos\theta$};
        \draw[->, line width=1.4pt, myr!85!black] (M) -- ++(62:0.95) node[above right] {$V = P\sin\theta$};
        \node[above left, font=\footnotesize] at (4.55,0.68) {$A_\theta$};
        \node[below, font=\footnotesize] at (2.0,-0.9) {Sección recta $A_0$};
    \end{tikzpicture}
    \caption{Descomposición de una fuerza axial $P$ sobre una sección inclinada $A_\theta$ en sus componentes normal $F$ y cortante $V$.}
    \label{fig:t8_seccion_inclinada}
\end{figure}

\nt{En la Figura~\ref{fig:t8_seccion_inclinada}, si el ángulo del plano inclinado respecto a la normal es $\theta$, la normal exterior al plano forma un ángulo $\theta$ con el eje de la barra, mientras que el plano de corte forma un ángulo $\theta$ con el plano vertical $A_0$. Por geometría básica, la relación entre áreas es $\cos\theta = \dfrac{A_0}{A_\theta} \implies A_\theta = \dfrac{A_0}{\cos\theta}$.}

Sustituyendo $F$, $V$ y $A_\theta = A_0/\cos\theta$, se obtienen las tensiones normal y tangencial en el plano inclinado:
\begin{align}
    \sigma(\theta) &= \frac{F}{A_\theta} = \frac{P \cos\theta}{A_0 / \cos\theta} = \frac{P}{A_0}\cos^2\theta \label{eq:t8_sigma_theta}\\
    \tau(\theta)   &= \frac{V}{A_\theta} = \frac{P \sin\theta}{A_0 / \cos\theta} = \frac{P}{A_0}\sin\theta\cos\theta = \frac{P}{2A_0}\sin(2\theta) \label{eq:t8_tau_theta}
\end{align}

De las ecuaciones~\eqref{eq:t8_sigma_theta} y~\eqref{eq:t8_tau_theta} se deducen dos casos límite fundamentales:
\begin{itemize}
    \item \textbf{Sección transversal recta ($\theta = 0^\circ$):} La tensión cortante es nula ($\tau = 0$) y la tensión normal alcanza su valor máximo absoluto:
          \begin{equation*}
              \sigma_{\text{máx}} = \frac{P}{A_0}
          \end{equation*}
    \item \textbf{Sección a $45^\circ$ ($\theta = 45^\circ$, donde $\sin\theta = \cos\theta = \frac{\sqrt{2}}{2}$):} La tensión cortante alcanza su valor máximo absoluto, igualándose en magnitud a la tensión normal sobre ese plano:
          \begin{equation*}
              \tau_{\text{máx}} = \sigma(45^\circ) = \frac{P}{2A_0}
          \end{equation*}
\end{itemize}

%=============================================================================
\section{Diseño de Uniones Atornilladas y Remachadas}

Muchas estructuras y recipientes de proceso están formados por chapas o perfiles que deben unirse entre sí. En materiales metálicos los medios de unión más habituales son los \textbf{tornillos}, los \textbf{remaches (roblones)} y las \textbf{soldaduras}.

\begin{itemize}
    \item \textbf{Uniones atornilladas:} Se ejecutan mediante tornillos (espiga roscada y cabeza), tuerca y arandela introducidos en agujeros practicados sobre las piezas a unir. Sus dimensiones geométricas y calidades están normalizadas (en España, tradicionalmente mediante la norma MV-106-1968 y posteriormente el Código Técnico y Eurocódigo 3). Permiten un montaje y desmontaje sencillo para mantenimiento.
    \item \textbf{Uniones remachadas o roblonadas:} Se realizan mediante roblones formados por una espiga cilíndrica lisa y una \textbf{cabeza de asiento}. El roblón se calienta a unos $1000\,^\circ\text{C}$, se introduce en caliente a través de los taladros alineados de las chapas y mediante golpeo o prensado del extremo libre de la espiga se conforma una segunda cabeza (cabeza de cierre), asegurando una unión permanente y muy ajustada al enfriar.
\end{itemize}

\subsection{Tipologías constructivas y parámetros geométricos}
Las uniones atornilladas y remachadas trabajan fundamentalmente a cortadura. En función de la línea de acción de la carga pueden estar \textbf{centradas} respecto al centroide del grupo de conectores o sometidas a \textbf{cargas excéntricas}. Atendiendo a la disposición de las chapas, se distinguen dos grandes familias (Figura~\ref{fig:t8_tipos_uniones}):
\begin{enumerate}
    \item \textbf{Uniones por solape:} Una chapa se monta directamente sobre la otra.
    \item \textbf{Uniones a tope:} Las dos piezas principales se enfrentan en el mismo plano y se enlazan mediante uno o dos \textbf{cubrejuntas} (en diseño industrial se emplea doble cubrejuntas para evitar excentricidades fuera del plano y duplicar los planos de corte).
\end{enumerate}

\begin{figure}[H]
    \centering
    \begin{tikzpicture}[>=Latex, thick, scale=0.95]
        % Solape 1 fila (corte lateral)
        \node[font=\bfseries\small] at (2.2, 2.3) {Unión por solape (1 fila)};
        \draw[fill=gray!18] (0.5, 1.35) rectangle (3.0, 1.65);
        \draw[fill=gray!28] (1.6, 1.05) rectangle (4.1, 1.35);
        \filldraw[fill=black!80] (2.2, 0.90) rectangle (2.4, 1.80);
        \filldraw[black] (2.3, 1.80) ellipse (0.18 and 0.08);
        \filldraw[black] (2.3, 0.90) ellipse (0.18 and 0.08);
        \draw[->, blue!70!black] (0.4, 1.50) -- (-0.2, 1.50) node[left] {$P$};
        \draw[->, blue!70!black] (4.2, 1.20) -- (4.8, 1.20) node[right] {$P$};
        \node[font=\footnotesize] at (0.8, 1.82) {$t_2$};
        \node[font=\footnotesize] at (3.8, 0.85) {$t_1$};

        % A tope 1 fila doble cubrejuntas (corte lateral)
        \node[font=\bfseries\small] at (8.7, 2.3) {Unión a tope con doble cubrejuntas (1 fila)};
        \draw[fill=gray!18] (6.2, 1.20) rectangle (8.65, 1.50);
        \draw[fill=gray!18] (8.75, 1.20) rectangle (11.2, 1.50);
        % Cubrejuntas superior e inferior
        \draw[fill=blue!12] (7.2, 1.50) rectangle (10.2, 1.72);
        \draw[fill=blue!12] (7.2, 0.98) rectangle (10.2, 1.20);
        % Dos roblones (uno a cada lado)
        \foreach \rx in {7.9, 9.5}{
            \filldraw[fill=black!80] (\rx-0.09, 0.85) rectangle (\rx+0.09, 1.85);
            \filldraw[black] (\rx, 1.85) ellipse (0.16 and 0.07);
            \filldraw[black] (\rx, 0.85) ellipse (0.16 and 0.07);
        }
        \draw[->, blue!70!black] (6.1, 1.35) -- (5.5, 1.35) node[left] {$P$};
        \draw[->, blue!70!black] (11.3, 1.35) -- (11.9, 1.35) node[right] {$P$};
    \end{tikzpicture}
    \caption{Esquema en sección de una unión por solape (corte simple) y una unión a tope con doble cubrejuntas (corte doble).}
    \label{fig:t8_tipos_uniones}
\end{figure}

Los parámetros geométricos característicos de una unión remachada o atornillada son:
\begin{itemize}
    \item \textbf{Diámetro del roblón ($d$):} Diámetro nominal de la espiga una vez remachada, coincidente con el diámetro del taladro practicado en la chapa.
    \item \textbf{Paso ($p$):} Distancia entre los centros de dos remaches consecutivos de una misma fila.
    \item \textbf{Paso entre filas ($p_t$):} Distancia transversal entre las líneas de centros de dos filas consecutivas.
    \item \textbf{Paso diagonal en unión al tresbolillo ($p_d$):} Distancia entre centros de remaches alternados de dos filas adyacentes dispuestas al tresbolillo.
    \item \textbf{Módulo repetitivo:} Porción mínima elemental de la unión (conjunto de remaches y ancho de chapa $p$) que se repite periódicamente a lo largo de la costura.
    \item \textbf{Espesores de las chapas ($t_1, t_2$):} En los desarrollos se toma por convenio $t_1 < t_2$, por lo que la chapa de espesor $t_1$ es la limitante en aplastamiento y tracción.
    \item \textbf{Número de módulos repetitivos ($n$):} Coincide en general con el número de roblones por fila.
\end{itemize}

\subsection{Hipótesis de diseño y tensiones admisibles}
Para el cálculo analítico de uniones remachadas o atornilladas se asumen las siguientes hipótesis simplificativas:
\begin{enumerate}
    \item El rozamiento existente entre las superficies de las chapas es despreciable.
    \item Los remaches llenan completamente los agujeros practicados.
    \item El reparto de esfuerzos entre conectores es uniforme: cada área resistente toma una fracción proporcional de la carga total.
    \item La resistencia al corte se evalúa mediante la tensión cortante media uniforme en la sección transversal del roblón.
\end{enumerate}

De los materiales que conforman los conectores y las placas se requieren tres tensiones admisibles:
\begin{itemize}
    \item Tensión cortante admisible del roblón o tornillo: $\tau_{\text{adm}}$.
    \item Tensión admisible a tracción de las chapas: $\sigma_{\text{trac,adm}}$.
    \item Tensión admisible a compresión o aplastamiento de las chapas: $\sigma_{\text{comp,adm}}$.
\end{itemize}

Toda unión trabajando bajo una carga total $P$ repartida entre $n$ módulos repetitivos ($P_M = P/n$) puede fallar por cuatro mecanismos diferenciados:
\begin{enumerate}[label=\alph*)]
    \item \textbf{Fallo por cortadura} de los vástagos de los remaches.
    \item \textbf{Fallo por compresión o aplastamiento} de la chapa más delgada tras el remache.
    \item \textbf{Fallo por tracción} de la sección neta de chapa comprendida entre agujeros consecutivos.
    \item \textbf{Fallo por desgarramiento marginal} de la chapa en la zona posterior entre el agujero y el borde libre de la placa.
\end{enumerate}

%-----------------------------------------------------------------------------
\subsection{Unión solapada con una fila de remaches}

Sea una unión por solape con 1 fila de remaches y $n$ módulos repetitivos sometida a una carga centrada total $P$. Cada módulo soporta $P_M = P/n$.

\paragraph{a) Fallo del roblón por cortadura ($R_1$).}
Cada módulo presenta 1 remache con 1 sola sección transversal de corte ($\frac{\pi}{4}d^2$). La tensión cortante actuante y la resistencia máxima por módulo en cortadura son:
\begin{equation}
    \tau = \frac{R_1}{\frac{\pi}{4}d^2} \quad \Longrightarrow \quad R_1 = \tau_{\text{adm}} \cdot 1\text{ fila} \cdot \frac{\pi}{4}d^2
    \label{eq:t8_R1_solape1}
\end{equation}
La capacidad total del sistema de $n$ módulos frente a cortadura es $n \cdot R_1$, y el número mínimo de módulos (o remaches) necesarios para resistir $P$ resulta:
\begin{equation}
    n_{\text{mín}} = \frac{P}{\tau_{\text{adm}} \cdot \frac{\pi}{4}d^2}
\end{equation}

\paragraph{b) Fallo de la chapa por compresión o aplastamiento ($R_2$).}
La fuerza se transmite del vástago a la chapa de menor espesor ($t_1$) proyectada sobre un rectángulo equivalente de área $t_1 \cdot d$:
\begin{equation}
    \sigma_{\text{comp}} = \frac{R_2}{t_1 \cdot d} \quad \Longrightarrow \quad R_2 = \sigma_{\text{comp,adm}} \cdot 1\text{ fila} \cdot t_1 \cdot d
    \label{eq:t8_R2_solape1}
\end{equation}
El sistema completo resiste $n \cdot R_2$, siendo el número mínimo de módulos frente a aplastamiento:
\begin{equation}
    n_{\text{mín}} = \frac{P}{\sigma_{\text{comp,adm}} \cdot t_1 \cdot d}
\end{equation}

\paragraph{c) Rotura por tracción entre remaches ($R_3$).}
En un ancho de módulo igual al paso $p$, la sección transversal neta de la chapa delgada descontando el taladro es $(p - d) \cdot t_1$:
\begin{equation}
    \sigma_{\text{trac}} = \frac{R_3}{(p - d)\cdot t_1} \quad \Longrightarrow \quad R_3 = \sigma_{\text{trac,adm}} \cdot (p - d) \cdot t_1
    \label{eq:t8_R3_general}
\end{equation}
Por tanto, la resistencia global a tracción de la chapa perforada es $n \cdot R_3$, y el número mínimo de módulos para resistir $P$ vale:
\begin{equation}
    n_{\text{mín}} = \frac{P}{\sigma_{\text{trac,adm}} \cdot (p - d) \cdot t_1}
\end{equation}

\paragraph{d) Rotura por desgarramiento del borde de la chapa.}
Para evitar que el remache arranque o desgarre la lengüeta final de chapa situada detrás del agujero, las normas de diseño establecen como distancia mínima al borde libre:
\begin{equation}
    D \ge 2\,d
    \label{eq:t8_distancia_borde}
\end{equation}

\begin{ImportanteBox}[title={Procedimiento General de Diseño Óptimo y Rendimiento de la Unión ($\mu$)}]
    Conocidos el diámetro del roblón $d$, los espesores $t_1, t_2$ y las tensiones admisibles ($\tau_{\text{adm}}, \sigma_{\text{trac,adm}}, \sigma_{\text{comp,adm}}$):
    \begin{enumerate}
        \item Evaluar la resistencia de cada módulo repetitivo a cortadura ($R_1$) y a aplastamiento ($R_2$).
        \item Igualar la menor de las dos anteriores a la resistencia de la chapa a tracción ($R_3$):
              \begin{equation*}
                  \text{mín}(R_1, R_2) = R_3 = \sigma_{\text{trac,adm}} \cdot (p - d) \cdot t_1
              \end{equation*}
              De esta igualdad se despeja el \textbf{paso óptimo $p$} entre roblones de una misma fila, equilibrando los mecanismos de fallo.
        \item Calcular el \textbf{rendimiento o eficacia de la unión ($\mu$)}, definido como la relación porcentual entre la resistencia de la chapa taladrada y la de la chapa continua sin agujeros:
              \begin{equation}
                  \mu = \frac{\sigma_{\text{trac,adm}} \cdot (p - d)\cdot t_1}{\sigma_{\text{trac,adm}} \cdot p \cdot t_1}\cdot 100 = \frac{p - d}{p}\cdot 100
                  \label{eq:t8_rendimiento_union}
              \end{equation}
    \end{enumerate}
\end{ImportanteBox}

%-----------------------------------------------------------------------------
\subsection{Unión solapada con dos filas de remaches enfrentados}
Cada módulo repetitivo engloba \textbf{dos roblones} dispuestos en dos filas paralelas, de modo que para $n$ módulos habrá un total de $2n$ remaches:
\begin{itemize}
    \item \textbf{Cortadura:} La carga $R_1$ del módulo se reparte entre los 2 roblones en corte simple:
          \begin{equation*}
              R_1 = (2\text{ filas}) \cdot \tau_{\text{adm}} \cdot \frac{\pi}{4}d^2
          \end{equation*}
    \item \textbf{Aplastamiento:} Cada módulo cuenta con 2 áreas proyectadas de aplastamiento en la chapa delgada $t_1$:
          \begin{equation*}
              R_2 = (2\text{ filas}) \cdot \sigma_{\text{comp,adm}} \cdot t_1 \cdot d
          \end{equation*}
    \item \textbf{Tracción en la chapa:} La sección neta crítica por módulo sigue siendo $(p - d)t_1$ a lo largo de la fila exterior, por lo que:
          \begin{equation*}
              R_3 = \sigma_{\text{trac,adm}} \cdot (p - d) \cdot t_1
          \end{equation*}
\end{itemize}
\nt{Si la unión tuviera $k$ filas de remaches enfrentados, basta multiplicar las resistencias $R_1$ y $R_2$ del módulo por $k$ filas, manteniéndose invariable la expresión de $R_3$ para la fila más solicitada.}

%-----------------------------------------------------------------------------
\subsection{Unión a tope con una fila de remaches}
Al enfrentar las dos chapas principales con \textbf{doble cubrejuntas} (uno arriba y otro abajo, dimensionados para no ser limitantes frente a la chapa principal de espesor $t_1$), cada lado de la junta transmite la totalidad de la carga $P$, y \textbf{cada roblón ofrece dos secciones transversales de corte} (cortadura doble):
\begin{itemize}
    \item \textbf{Cortadura doble:} Cada remache aporta $2$ planos resistentes:
          \begin{equation*}
              R_1 = (1\text{ fila}) \cdot 2 \cdot \tau_{\text{adm}} \cdot \frac{\pi}{4}d^2
          \end{equation*}
    \item \textbf{Aplastamiento en chapa principal:} La carga se apoya en una sola chapa principal de espesor $t_1$:
          \begin{equation*}
              R_2 = (1\text{ fila}) \cdot \sigma_{\text{comp,adm}} \cdot t_1 \cdot d
          \end{equation*}
    \item \textbf{Tracción en chapa principal:} $R_3 = \sigma_{\text{trac,adm}} \cdot (p - d) \cdot t_1$.
\end{itemize}

%-----------------------------------------------------------------------------
\subsection{Unión a tope con dos filas de remaches}
Cada módulo repetitivo posee 2 roblones por lado de chapa principal, trabajando cada uno de ellos en cortadura doble gracias al doble cubrejuntas:
\begin{itemize}
    \item \textbf{Cortadura doble con 2 filas:} $R_1 = (2\text{ filas}) \cdot 2 \cdot \tau_{\text{adm}} \cdot \frac{\pi}{4}d^2 = 4\,\tau_{\text{adm}}\frac{\pi}{4}d^2$.
    \item \textbf{Aplastamiento:} $R_2 = (2\text{ filas}) \cdot \sigma_{\text{comp,adm}} \cdot t_1 \cdot d$.
    \item \textbf{Tracción:} $R_3 = \sigma_{\text{trac,adm}} \cdot (p - d) \cdot t_1$.
\end{itemize}

\begin{table}[H]
    \centering
    \caption{Resumen comparativo de resistencias por módulo repetitivo en uniones remachadas o atornilladas centradas (para $n$ módulos, $P_{\text{total}} = n\cdot \text{mín}(R_1,R_2,R_3)$).}
    \adjustbox{max width=\textwidth}{%
        \renewcommand{\arraystretch}{1.28}
        \setlength{\tabcolsep}{7pt}
        \small
        \begin{tabular}{@{}lcccc@{}}
            \toprule
            \textbf{Tipo de Unión} & \textbf{\shortstack{Cortadura\\($R_1$)}} & \textbf{\shortstack{Aplastamiento\\($R_2$)}} & \textbf{\shortstack{Tracción neta\\($R_3$)}} & \textbf{\shortstack{Roblones totales\\en la junta}} \\
            \midrule
            Solape (1 fila) & $\tau_{\text{adm}}\,\frac{\pi}{4}d^2$ & $\sigma_{\text{comp,adm}}\,t_1 d$ & $\sigma_{\text{trac,adm}}(p-d)t_1$ & $n$ \\
            Solape (2 filas enfrentadas) & $2\,\tau_{\text{adm}}\,\frac{\pi}{4}d^2$ & $2\,\sigma_{\text{comp,adm}}\,t_1 d$ & $\sigma_{\text{trac,adm}}(p-d)t_1$ & $2n$ \\
            A tope doble cubrejuntas (1 fila) & $2\,\tau_{\text{adm}}\,\frac{\pi}{4}d^2$ & $\sigma_{\text{comp,adm}}\,t_1 d$ & $\sigma_{\text{trac,adm}}(p-d)t_1$ & $2n$ ($n$ por lado) \\
            A tope doble cubrejuntas (2 filas) & $4\,\tau_{\text{adm}}\,\frac{\pi}{4}d^2$ & $2\,\sigma_{\text{comp,adm}}\,t_1 d$ & $\sigma_{\text{trac,adm}}(p-d)t_1$ & $4n$ ($2n$ por lado) \\
            \bottomrule
        \end{tabular}%
    }
    \label{tab:t8_resumen_uniones}
\end{table}

\begin{table}[H]
    \centering
    \caption{Esfuerzos permisibles de diseño para conectores de acero estructural (remaches ASTM A502 y tornillos ASTM A325 / A490).}
    \adjustbox{max width=0.88\textwidth}{%
        \renewcommand{\arraystretch}{1.24}
        \setlength{\tabcolsep}{9pt}
        \small
        \begin{tabular}{@{}llcc@{}}
            \toprule
            \textbf{Tipo de Conector} & \textbf{Designación / Grado} & \textbf{\shortstack{Esfuerzo cortante\\permisible $\tau_{\text{adm}}$ (MPa)}} & \textbf{\shortstack{Esfuerzo de tracción\\permisible $\sigma_{\text{trac,adm}}$ (MPa)}} \\
            \midrule
            \multirow{2}{*}{\textbf{Remaches}} & ASTM A502 Grado 1 & 121 & 159 \\
                                               & ASTM A502 Grado 2 & 152 & 200 \\
            \midrule
            \multirow{2}{*}{\textbf{Tornillos}} & ASTM A325 & 121 & 303 \\
                                                & ASTM A490 & 152 & 372 \\
            \bottomrule
        \end{tabular}%
    }
    \label{tab:t8_esfuerzos_astm}
\end{table}
